% Proof portfolio template for Math 250, Meredith College
% Phillip Andreae
% October 16, 2019

% Students should also consult Overleaf resources and advice I will post on Brightspace.


%%%%%%%%%%%%%%%%%%%%%%%%%%%%%%%%
% This part is called the preamble.
%%%%%%%%%%%%%%%%%%%%%%%%%%%%%%%%
\documentclass[12pt]{article} % specifies font size and document type
\usepackage{amsmath, amssymb, amsthm} % these packages are necessary for math formatting

% A percent sign comments the rest of the line out,
% so these lines do not effect what is displayed in your document.

\theoremstyle{definition}
% Here I'm defining some "theorem environments":
\newtheorem*{defn}{Definition}
\newtheorem*{prop}{Proposition}
\newtheorem*{theorem}{Theorem}

% (end of the preamble)
%%%%%%%%%%%%%%%%%%%%%%%%%%%%%%%%




%%%%%%%%%%%%%%%%%%%%%%%%%%%%%%%%
% The body of the document starts after \begin{document}
%%%%%%%%%%%%%%%%%%%%%%%%%%%%%%%%
\begin{document}

\title{Math 250 Proof Portfolio}
\author{Your Name}
\date{December 4, 2019}

\maketitle % This displays the title, author, and date that were set above.



%%%%%%%%%%%%%%%%%%%%%%%%%%%%%%%%%%%%%%%%
% Definitions
%%%%%%%%%%%%%%%%%%%%%%%%%%%%%%%%%%%%%%%%
\section*{Definitions} % The * tells it not to number this section

% You should include here the definition
% of every term that appears in one of your proofs.
% Below are some sample  definitions.

\begin{defn}
An integer $x$ is \textit{even}
if $x = 2a$ for some integer $a$.
 % I'm putting every mathematical expression, including
 % variable names, between dollar signs.
\end{defn}

\begin{defn}
 A \textit{rectangle} is a quadrilateral
 all of whose angles are right angles.
A \textit{square} is a rectangle
all of whose sides are congruent.
\end{defn}




%%%%%%%%%%%%%%%%%%%%%%%%%%%%%%%%%%%%%%%%
% Proofs
%%%%%%%%%%%%%%%%%%%%%%%%%%%%%%%%%%%%%%%%

%%%%%%%%%%%%%%%%%%%%%%%%%%%%%%%%%%%%%%%%
\section{Proof by cases} % It will number the sections without *'s automatically

\begin{prop} % I'm calling this a proposition since it doesn't seem famous or important enough to call it a "theorem".
If $X$ is a square, then $X$ is a rectangle.
\end{prop}

\begin{proof}
 Suppose $X$ is a square.
 
 I'm going to demonstrate some of the cool ways
 that \LaTeX  can format mathematical expressions!
 Let $s$ be the side length of $X$
 and $A$ be the area of $X$.

 Then we have $A = s^2$.
 Now suppose that $s$ is an even integer.
 By definition, this means that $s = 2a$ for some $a \in \mathbb{Z}$.
 Then the area of $X$ is
 \begin{equation*}
     A = s^2 = (2a)^2 = 4a^2,
 \end{equation*}
 which we also could have formatted over multiple lines as
  \begin{align*}
     A &= s^2 \\
     &= (2a)^2 \\
     &= 4a^2,
 \end{align*}
 which shows that $A$ is a multiple of $4$,
 i.e., $A \equiv 0 \pmod{4}$.
 
 We can format set-builder notation like this:
 \begin{equation*}
     \{ 2n : n \in \mathbb{Z} \} \subseteq \{ m \in \mathbb{Z} : 4|m \}.
 \end{equation*}
 
 Now suppose that the side length of $X$
 is a function of $t$, where $t \in \mathbb{R}$.
 Then by the chain rule, the derivative of
 the area is $\frac{dA}{dt} = 2s \frac{ds}{dt}$,
 which we could also have formatted as
 $\displaystyle \frac{dA}{dt} = 2s \frac{ds}{dt}$
 if we wanted to make the fractions larger.
 The net change in the area from $t=t_1$
 to $t=t_2$ is
 \begin{equation*}
     A(t_1) - A(t_2) = \int_{t_1}^{t_2} \frac{dA}{dt} \, dt.
 \end{equation*}
 Notice how I put a period at the end of that sentence!
 
 I said this was a proof by cases, so let's see how to format lists.
Either $X$ is small or $X$ is large. Let's consider those two cases separately:
\begin{enumerate}
    \item Suppose $X$ is small \dots
    \item Suppose $X$ is large \dots
\end{enumerate}
Here's how to create a bulleted list:
\begin{itemize}
    \item \textit{Case 1:} Suppose $X$ is small \dots
    \item \textit{Case 2:} Suppose $X$ is large \dots
\end{itemize}
 
 By the way, since a square is by definition
 a rectangle all of whose sides are congruent,
 $X$ is also a rectangle.
 This completes the proof.
 Marvel at how \LaTeX will create the ``end-of-proof-box"
 automatically:
\end{proof}


%%%%%%%%%%%%%%%%%%%%%%%%%%%%%%%%%%%%%%%%
\section{Congruence modulo $n$}


%%%%%%%%%%%%%%%%%%%%%%%%%%%%%%%%%%%%%%%%
\section{Proof by contrapositive}


%%%%%%%%%%%%%%%%%%%%%%%%%%%%%%%%%%%%%%%%
\section{Proof by contradiction}


%%%%%%%%%%%%%%%%%%%%%%%%%%%%%%%%%%%%%%%%
\section{Induction}


%%%%%%%%%%%%%%%%%%%%%%%%%%%%%%%%%%%%%%%%
\section{A proof that a function is bijective}


%%%%%%%%%%%%%%%%%%%%%%%%%%%%%%%%%%%%%%%%
\section{The triangle inequality}


%%%%%%%%%%%%%%%%%%%%%%%%%%%%%%%%%%%%%%%%
\section{The Pythagorean theorem}

Use brackets to give a theorem a name in parentheses:
\begin{theorem}[Pythagorean theorem] 
For every right triangle with legs of lengths $a$ and $b$
		and hypotenuse of length $c$,
		we have that $a^2 + b^2 = c^2$.
\end{theorem}

Compare to:
\begin{theorem}
This theorem does not have a name
in parentheses since I did not use the brackets.
\end{theorem}

You could cite a source by saying something like:
The idea of the following proof, which is originally due to Euclid,
has been borrowed from \cite{pyth}. % This cites the source I named "pyth" in my bibliography

\begin{proof}
    Here's the proof.
\end{proof}


%%%%%%%%%%%%%%%%%%%%%%%%%%%%%%%%%%%%%%%%
\section{A proof not from this course}



%%%%%%%%%%%%%%%%%%%%%%%%%%%%%%
% Bibliography
%%%%%%%%%%%%%%%%%%%%%%%%%%%%%%
\begin{thebibliography}{1}

\bibitem{hammack} R. Hammack, \textit{Book of Proof}.
  Third edition. 2018.
  % I've given this source the name "hammack",
  % so I can use the command \cite{hammack} when I want to cite it.

\bibitem{pyth} Wikipedia contributors,
\textit{Wikipedia},
``Pythagorean theorem".
https://en.wikipedia.org/w/index.php?title=
Pythagorean\_theorem  \\
\&oldid=918079404.
Online; accessed October 15, 2019.


  \end{thebibliography}


\end{document}